
Soit alors \(X=\underline{\underline{D}}(\underline{E})\times_S \underline{\underline{D}}(\underline{F})\). Le \(\lambda\)-homomorphisme \(K^{\bullet}(S)\to K^{\bullet}(X)\) est injectif d'après 3.1. Par conséquent, on peut, pour montrer les relations V 2.4.2, se placer dans \(K^{\bullet}(X)\). Comme dans \(K^{\bullet}(X)\) les éléments \(E\) et \(F\) se décomposent en somme de classe de faisceaux inversibles, il suffit, par additivité, de montrer ces mêmes relations pour des éléments \(\xi,\eta\), classes de faisceaux inversibles ; elles peuvent encore s'écrire~:
\[
  1+\xi\eta t=(1+\xi t)\circ(1+\eta t);
  \qquad
  \lambda_u(1+\xi t)=1+\{1+\xi t\}u,
\]
en les mettant sous la forme V 2.4.1. Elles ne sont alors autres que les relations V 2.3.1 et 2.3.3 qui définissent la structure de \(\lambda\)-anneau de \(1+\widehat{K^{\bullet}(S)[[t]]}^{+}\).

Indiquons le résultat plus général suivant~:

\textbf{Théorème 3.3} \uline{Soit} \((T,\underline{A})\) \uline{un topos commutativement annelé}; \uline{on désigne par} \(K\) \uline{le groupe de Grothendieck des} \(A\)-\uline{Modules localement libres} de rang borné sur \(T\). \uline{Alors}, \(K\), \uline{muni de sa pré}-\(\lambda\)-\uline{structure canonique}~(V 2.2 b)), \uline{est} \uline{un} \(\lambda\)-\uline{anneau}.

Pour montrer que \(K\) est un \(\lambda\)-anneau, il faut montrer les relations V 2.4.1, et par additivité, il suffit de les montrer pour des éléments \(E,F\) qui sont classes de \(\underline{A}\)-Modules localement libres de rang borné. Comme les rangs de \(\underline{E}\) et \(\underline{F}\) sont bornés, on peut décomposer \(T\) en «~somme disjointe~» finie d'ouverts sur lesquels \(\underline{E}\) et \(\underline{F}\) sont de rang constant. En utilisant la commutation du \(K^{\bullet}\) aux sommes directes finies, et en mettant les relations V 2.4.1 sous la forme V 2.4.2, on voit qu'il suffit de montrer les relations V 2.4.2, dans chacun des \(K^{\bullet}\) des ouverts considérés, c'est-à-dire qu'il suffit de montrer les relations V 2.4.2 pour des éléments \(E,F\), classes de \(\underline{A}\)-Modules \(\underline{E},\underline{F}\) localement libres de rang constant.

